\section{Levantamento de Requisitos}

Em coerência com a Engenharia de Requisitos discutida na Base
Teórica, o levantamento apresentado nesta seção corresponde
às etapas de \textit{Análise} do ciclo de BPM: a partir do
diagnóstico do processo atual, busca-se identificar os
gargalos operacionais e transformá-los em requisitos capazes
de orientar o redesenho do processo, apresentado na seção
seguinte. O levantamento de requisitos foi realizado a partir
das respostas do formulário estruturado aplicado aos três
participantes definidos na metodologia: uma Analista de
Qualidade e Clientes, uma Supervisora de Qualidade e um
Gerente da Área. Os três participantes registraram
concordância voluntária com a pesquisa. As perguntas foram
segmentadas conforme o nível de atuação, de modo que os
campos sem resposta nos blocos específicos não representam
omissão, mas sim questões direcionadas a outro perfil
organizacional.

A análise das respostas foi conduzida por agrupamento
temático, considerando cinco eixos: processo atual,
dificuldades e riscos, regras de negócio, funcionalidades
esperadas e benefícios gerenciais. A partir desses eixos,
as necessidades relatadas foram convertidas em requisitos
funcionais, requisitos não funcionais e regras de negócio.
Para a priorização foi utilizada a técnica MoSCoW,
considerando como \textit{Must Have} os requisitos
necessários para que o fluxo opere com rastreabilidade e
governança mínimas, como \textit{Should Have} aqueles que
aumentam a qualidade operacional sem impedir a operação
básica, e como \textit{Could Have} as funcionalidades de
evolução que podem ser incorporadas em ciclos posteriores.

\subsection{Diagnóstico do Processo Atual}

A resposta da Analista de Qualidade e Clientes permitiu
caracterizar o processo operacional utilizado antes da
solução proposta. Os chamados internos eram registrados no
SAP Service, ferramenta também utilizada como CRM para
atendimentos e reclamações de clientes externos. Para
viabilizar esse registro, as ocorrências internas eram
associadas a um cadastro denominado ``Experiência do
Cliente'', no qual era criado um ticket e realizado o
direcionamento à área responsável. Paralelamente ao registro
no sistema, a chegada do chamado era comunicada manualmente
aos responsáveis por meio de um grupo de WhatsApp.

Esse modelo apresentava duas fragilidades centrais. A
primeira estava relacionada à mistura, em um mesmo ambiente,
de chamados internos e manifestações externas de clientes,
o que poderia distorcer a leitura de volumes e dificultar a
governança das ocorrências internas. A segunda estava
relacionada à dependência de atividades manuais para
comunicação e acompanhamento, especialmente quanto à
obtenção de retorno das áreas responsáveis.

A Analista também diferenciou os dois tipos de chamado
utilizados no processo. O chamado interno corresponde a uma
problemática organizacional que necessita de análise e
tratativa por outra área. O chamado massivo corresponde a
uma ocorrência de mesmo problema capaz de impactar três ou
mais clientes, normalmente associada a falhas técnicas ou
sistêmicas com maior potencial de impacto. Essa distinção
passou a ser considerada na modelagem do processo e nas
regras de priorização.

Na perspectiva da Supervisora de Qualidade, a proposta de
um ambiente específico para registro e acompanhamento foi
considerada aderente às necessidades operacionais. A
respondente destacou como benefícios esperados a maior
agilidade no direcionamento, a rastreabilidade das
tratativas, a simplificação do processo em comparação ao
SAP Service, a notificação imediata da área responsável e a
possibilidade de consolidação das informações para análise
gerencial. Como pontos de atenção, foram identificadas a
necessidade de registrar atualizações ou justificativas
quando a solução depender de fatores externos e a
necessidade de permitir a transferência do chamado quando o
responsável inicialmente definido não for o correto.

Na perspectiva gerencial, foram confirmadas como principais
fricções a dificuldade de governança dos chamados internos,
a falta de padronização dos registros e do direcionamento e
a dificuldade das áreas executoras em visualizar suas
pendências. O Gerente da Área avaliou que o fluxo proposto
possui aderência aos objetivos estratégicos da gerência e
da organização, destacando como benefícios esperados a
rastreabilidade, a padronização, a disponibilidade de
indicadores gerenciais e a redução do tempo de resolução.
O fluxo foi aprovado, em princípio, sem ressalvas.

A Tabela~\ref{tab:sintese-levantamento} sintetiza as
contribuições de cada perfil para a especificação da
solução.

\begin{table}[H]
\caption{Síntese do levantamento de requisitos por perfil}
\label{tab:sintese-levantamento}
\centering
\begin{tabular}{p{2.4cm} p{5cm} p{5cm}}
\hline
\textbf{Perfil} & \textbf{Principais evidências} &
\textbf{Implicações para a solução} \\
\hline
Analista de Qualidade e Clientes &
Uso do SAP Service para ocorrências internas e externas;
comunicação complementar por WhatsApp; dificuldade de
retorno; necessidade de protocolo, evidências, regras,
causa raiz, solução, SLA e indicadores. &
Criar canal próprio; padronizar a abertura; gerar protocolo;
automatizar notificações; registrar causa raiz e ação;
separar chamados internos de manifestações externas. \\
\hline
Supervisora de Qualidade &
Necessidade de classificação e solução por prazo;
rastreabilidade; transferência de responsável; registro de
atualizações e justificativas quando o SLA não puder ser
cumprido. &
Definir papéis de classificação; controlar mudança de
responsável; manter histórico; prever mecanismo de
justificativa e acompanhamento; estruturar regras de SLA. \\
\hline
Gerente da Área &
Dificuldade de governança, padronização e acesso às
pendências; necessidade de indicadores; aprovação do fluxo
proposto. &
Centralizar os registros; disponibilizar visão gerencial;
priorizar volume, tempo médio de resolução e distribuição
por área; garantir aderência ao processo organizacional. \\
\hline
\end{tabular}
\raggedright
\par
{\fontsize{10}{12}\selectfont Fonte: Elaborado pelo autor
com base na pesquisa de campo}
\end{table}

\subsection{Requisitos Funcionais}

Os requisitos funcionais representam os serviços que a
solução deve disponibilizar aos usuários e às áreas
envolvidas no processo. A Tabela~\ref{tab:requisitos-funcionais}
apresenta os requisitos consolidados a partir das respostas
dos participantes.

\begin{table}[H]
\caption{Requisitos funcionais da solução}
\label{tab:requisitos-funcionais}
\centering
\begin{tabular}{p{1.0cm} p{8.7cm} p{2.3cm}}
\hline
\textbf{ID} & \textbf{Requisito funcional} &
\textbf{Prioridade} \\
\hline
RF01 & Disponibilizar um canal único e específico para
abertura de chamados internos, separado do ambiente de
reclamações externas de clientes. & Must Have \\
\hline
RF02 & Permitir que usuários autorizados de diferentes
áreas registrem novos chamados internos. & Must Have \\
\hline
RF03 & Permitir a identificação do tipo de chamado,
diferenciando, no mínimo, chamados Internos e Massivos. &
Must Have \\
\hline
RF04 & Registrar de forma padronizada as informações
necessárias à abertura, incluindo empresa, tipo, assunto,
descrição do problema e evidências, além de dados
operacionais aplicáveis ao caso. & Must Have \\
\hline
RF05 & Gerar automaticamente um número único de protocolo
para cada chamado registrado. & Must Have \\
\hline
RF06 & Confirmar a abertura ao solicitante e informar o
protocolo e o link de acesso ao registro. & Must Have \\
\hline
RF07 & Notificar a equipe de Qualidade/SAC sobre novos
chamados que aguardam classificação e direcionamento. &
Must Have \\
\hline
RF08 & Permitir à Qualidade/SAC definir prioridade, área
responsável, tipologia, responsável pela tratativa e status
do chamado. & Must Have \\
\hline
RF09 & Notificar automaticamente o responsável quando um
chamado for atribuído para sua tratativa. & Must Have \\
\hline
RF10 & Permitir a transferência do chamado para outro
responsável, preservando a rastreabilidade da alteração. &
Must Have \\
\hline
RF11 & Permitir o registro de causa raiz e ação corretiva
pela área responsável pela tratativa. & Must Have \\
\hline
RF12 & Encaminhar o chamado tratado para análise da
Qualidade/SAC antes do encerramento definitivo. &
Must Have \\
\hline
RF13 & Permitir à Qualidade/SAC concluir ou cancelar o
chamado, registrando o motivo do cancelamento quando
aplicável. & Must Have \\
\hline
RF14 & Notificar o solicitante e os responsáveis quando o
chamado for concluído ou cancelado. & Must Have \\
\hline
RF15 & Controlar os prazos de classificação e solução de
acordo com o tipo e a prioridade do chamado, possibilitando
o acompanhamento do SLA. & Must Have \\
\hline
RF16 & Permitir o registro de comentários, atualizações ou
justificativas para chamados cuja tratativa dependa de
fatores externos ou ultrapasse o prazo inicialmente
definido. & Should Have \\
\hline
RF17 & Permitir o registro da quantidade de protocolos
relacionados quando a ocorrência afetar múltiplos clientes
ou registros. & Should Have \\
\hline
RF18 & Disponibilizar categorias ou tipologias reutilizáveis
para apoiar a classificação e a análise de recorrências. &
Should Have \\
\hline
RF19 & Disponibilizar indicadores gerenciais de volume,
tempo médio de resolução, distribuição por área, tipos de
problema, chamados tratados e cumprimento de SLA. &
Must Have \\
\hline
\end{tabular}
\raggedright
\par
{\fontsize{10}{12}\selectfont Fonte: Elaborado pelo autor
com base na pesquisa de campo}
\end{table}

\subsection{Requisitos Não Funcionais}

Além das funcionalidades, o levantamento evidenciou
características de qualidade necessárias para que a solução
seja incorporada à rotina organizacional. Esses requisitos
estão apresentados na Tabela~\ref{tab:requisitos-nao-funcionais}.

\begin{table}[H]
\caption{Requisitos não funcionais da solução}
\label{tab:requisitos-nao-funcionais}
\centering
\begin{tabular}{p{1.1cm} p{8.5cm} p{2.4cm}}
\hline
\textbf{ID} & \textbf{Requisito não funcional} &
\textbf{Prioridade} \\
\hline
RNF01 & A solução deve centralizar os registros em uma fonte
única e permitir rastreabilidade desde a abertura até o
encerramento. & Must Have \\
\hline
RNF02 & A interface de abertura deve ser simples, objetiva e
adequada ao uso em ambiente corporativo, reduzindo o número
de etapas necessárias para registrar uma ocorrência. &
Must Have \\
\hline
RNF03 & A solução deve utilizar, prioritariamente, os
recursos do ecossistema Microsoft 365 já adotado pela
organização, favorecendo integração e manutenção. &
Must Have \\
\hline
RNF04 & O acesso aos registros e às ações do processo deve
respeitar as permissões corporativas e os diferentes papéis
de solicitante, Qualidade/SAC e área responsável. &
Must Have \\
\hline
RNF05 & As alterações relevantes de status, responsável,
tratativa e encerramento devem permanecer auditáveis no
histórico do registro. & Must Have \\
\hline
RNF06 & As automações devem evitar notificações duplicadas e
reprocessamentos causados pelas próprias atualizações
realizadas no item. & Must Have \\
\hline
RNF07 & A interface incorporada ao SharePoint deve adaptar-se
ao espaço disponível na página, preservando legibilidade e
usabilidade. & Should Have \\
\hline
\end{tabular}
\raggedright
\par
{\fontsize{10}{12}\selectfont Fonte: Elaborado pelo autor}
\end{table}

\subsection{Regras de Negócio}

As respostas da pesquisa permitiram definir regras
operacionais que orientam o comportamento do fluxo e a
distribuição de responsabilidades entre os participantes.
A Tabela~\ref{tab:regras-negocio} apresenta as principais
regras identificadas.

\begin{table}[H]
\caption{Regras de negócio identificadas}
\label{tab:regras-negocio}
\centering
\begin{tabular}{p{1.0cm} p{10.9cm}}
\hline
\textbf{ID} & \textbf{Regra de negócio} \\
\hline
RN01 & Todo chamado deve ser registrado assim que a
problemática for identificada, garantindo rastreabilidade
desde o início do processo. \\
\hline
RN02 & A Qualidade/SAC é responsável pela classificação e
pelo direcionamento inicial do chamado para a área
responsável pela tratativa. \\
\hline
RN03 & Chamados Massivos devem possuir tratamento mais
restritivo em razão do potencial impacto sobre maior
quantidade de clientes e indicadores. \\
\hline
RN04 & Na pesquisa de campo, a classificação proposta para
Chamados Massivos foi Alta e para Chamados Internos foi
Média. \\
\hline
RN05 & A classificação inicial de um Chamado Massivo deve
ocorrer em até 30 minutos após a abertura; para Chamados
Internos, em até um dia após a abertura. \\
\hline
RN06 & A solução de um Chamado Massivo deve ocorrer em até
24 horas após a classificação; para Chamados Internos, em
até três dias após a classificação. \\
\hline
RN07 & A área responsável deve registrar, no mínimo, a causa
raiz identificada e a ação corretiva adotada antes de
submeter a tratativa para análise da Qualidade/SAC. \\
\hline
RN08 & A conclusão definitiva do chamado deve ser realizada
pela Qualidade/SAC após análise da tratativa registrada pela
área responsável. \\
\hline
RN09 & Quando houver transferência de responsável, a
alteração deve preservar a rastreabilidade do chamado e
notificar o novo responsável. \\
\hline
RN10 & Em situações extraordinárias de impacto imediato,
como indisponibilidade generalizada de sistemas, o processo
deve permitir tratamento prioritário sem impedir o registro
formal da ocorrência. \\
\hline
\end{tabular}
\raggedright
\par
{\fontsize{10}{12}\selectfont Fonte: Elaborado pelo autor
com base na pesquisa de campo}
\end{table}

\subsection{Decisões de Modelagem Derivadas dos Requisitos}

Durante a transformação dos requisitos em solução, algumas
decisões de modelagem foram adotadas para reduzir erros de
direcionamento e separar responsabilidades. Embora a
Analista tenha indicado ``Área responsável'' entre as
informações desejáveis na abertura, a mesma respondente
também solicitou que a área de Clientes/Qualidade atuasse
como filtro interno para identificar o destino correto.
A Supervisora, por sua vez, apontou a necessidade de
transferência quando o direcionamento inicial estivesse
incorreto. Em razão dessa convergência, a definição da área
responsável, da prioridade, da tipologia e do responsável
pela tratativa foi posicionada na etapa de classificação,
sob responsabilidade da Qualidade/SAC, e não no formulário
inicial do solicitante.

Também foi estabelecida uma etapa intermediária entre a
tratativa e o encerramento. Após o preenchimento da causa
raiz e da ação corretiva, o chamado deve retornar para
análise da Qualidade/SAC. Essa decisão atende à necessidade
de obrigatoriedade de posicionamento, análise de causa e
solução indicada no levantamento e impede que a própria área
responsável encerre diretamente um chamado sem validação.

Por fim, o levantamento demonstrou que a automação não deve
ser compreendida apenas como mecanismo de registro. A
necessidade de governança, rastreabilidade e indicadores
foi recorrente nos três níveis consultados, tornando a
centralização dos dados e a manutenção de um histórico
estruturado elementos centrais da solução.
